

In this project, we formed a corpus of documents related to the eucalyptus tree, which we then used to quantitatively analyze the discourse around the plant in the French-speaking world, between the European discovery of the tree and the end of the First World War. In this exploratory analysis, we tried to discover which actors shaped the discourse around the tree, and how they did it. Concerning the actors, the data shows several different types. We see that while at the beginning the discourse was enclosed within scientific publications, the press started taking control after 1870. While these documents are not as in-depth as the scientific monographs, they largely participate in raising general awareness about the plant. Moreover, by looking at places of publication, we see that there are different actors at play, between mainland France and its colonies. This can be seen in the way eucalyptus is talked about, between the cases of Lyon and Madagascar. The first driven by the generalist press and advertisements about the medicinal virtues of the plant, while the latter being used in the official state press, confined to institutional topics. This difference of treatment does portray the eucalyptus tree as a colonization tool. However, we do also find differences within mainland France, notably between the north, the south and Paris. While the south is mainly linked with adverts, we see that the north, where no plantations are found, the discourse is often slightly more scientific in nature (with most topics being linked to medicine, geography or botanic), with fewer generalist journals. Moreover, when generalist press appears, it seems to be linked with religious institutions, as seen with the Trappist's plantations in Roman Campagna. In Paris meanwhile, we find the overwhelmingly majority of monograph publication and thus the more scientific fields of medicine, while at the same time finding most of the press production of the country. It is, unsurprisingly, the cultural and scientific leader of the French-speaking world, and thus of the eucalyptus discourse.

However, while this quantitative approach works to understand who the actors are and how they tend to shape the discourse around eucalyptus, it struggles to understand their motivation. Indeed, this project only captures topics, without capturing the thought process behind certain political or editorial decisions, specific to each document. And while theories can be emitted, this approach is better suited as a complement to a proper qualitative research. In conclusion, this project contributes to the research on the discourse around eucalyptus, by showing the effects of certain policies in the general discourse, as well as finding new areas of potential research, while not having the ambition of replacing proper qualitative studies, key to better explaining the dynamics explored in this project.
